\documentclass[otchet]{SCWorks}
% Тип обучения (одно из значений):
%    bachelor   - бакалавриат (по умолчанию)
%    spec       - специальность
%    master     - магистратура
% Форма обучения (одно из значений):
%    och        - очное (по умолчанию)
%    zaoch      - заочное
% Тип работы (одно из значений):
%    coursework - курсовая работа (по умолчанию)
%    referat    - реферат
%  * otchet     - универсальный отчет
%  * nirjournal - журнал НИР
%  * digital    - итоговая работа для цифровой кафедры
%    diploma    - дипломная работа
%    pract      - отчет о научно-исследовательской работе
%    autoref    - автореферат выпускной работы
%    assignment - задание на выпускную квалификационную работу
%    review     - отзыв руководителя
%    critique   - рецензия на выпускную работу

% * Добавлены вручную. За вопросами к @mchernigin
\usepackage{preamble}

\begin{document}

% Кафедра (в родительном падеже)
\chair{математической кибернетики и компьютерных наук}

% Тема работы
\title{Анализ красно-чёрного дерева}

% Курс
\course{2}

% Группа
\group{211}

% Факультет (в родительном падеже) (по умолчанию "факультета КНиИТ")
% \department{факультета КНиИТ}

% Специальность/направление код - наименование
 \napravlenie{02.03.02 "--- Фундаментальная информатика и информационные технологии}
% \napravlenie{02.03.01 "--- Математическое обеспечение и администрирование информационных систем}
% \napravlenie{09.03.01 "--- Информатика и вычислительная техника}
%\napravlenie{09.03.04 "--- Программная инженерия}
% \napravlenie{10.05.01 "--- Компьютерная безопасность}

% Для студентки. Для работы студента следующая команда не нужна.
% \studenttitle{студентки}

% Фамилия, имя, отчество в родительном падеже
\author{Бурова Арсения Александровича}

% Заведующий кафедрой 
\chtitle{доцент, к.\,ф.-м.\,н.}
\chname{С.\,В.\,Миронов}

% Руководитель ДПП ПП для цифровой кафедры (перекрывает заведующего кафедры)
% \chpretitle{
%     заведующий кафедрой математических основ информатики и олимпиадного\\
%     программирования на базе МАОУ <<Ф"=Т лицей №1>>
% }
% \chtitle{г. Саратов, к.\,ф.-м.\,н., доцент}
% \chname{Кондратова\, Ю.\,Н.}

% Научный руководитель (для реферата преподаватель проверяющий работу)
\satitle{доцент, к.\,ф.-м.\,н.} %должность, степень, звание
\saname{Ю.\,Н.\,Кондратова}

% Руководитель практики от организации (руководитель для цифровой кафедры)
\patitle{доцент, к.\,ф.-м.\,н.}
\paname{С.\,В.\,Миронов}

% Руководитель НИР
\nirtitle{доцент, к.\,п.\,н.} % степень, звание
\nirname{В.\,А.\,Векслер}

% Семестр (только для практики, для остальных типов работ не используется)
\term{2}

% Наименование практики (только для практики, для остальных типов работ не
% используется)
\practtype{учебная}

% Продолжительность практики (количество недель) (только для практики, для
% остальных типов работ не используется)
\duration{2}

% Даты начала и окончания практики (только для практики, для остальных типов
% работ не используется)
\practStart{01.07.2024}
\practFinish{13.01.2024}

% Год выполнения отчета
\date{2026}

\maketitle

% Включение нумерации рисунков, формул и таблиц по разделам (по умолчанию -
% нумерация сквозная) (допускается оба вида нумерации)
\secNumbering

\tableofcontents

% Раздел "Обозначения и сокращения". Может отсутствовать в работе
% \abbreviations
% \begin{description}
%     \item ... "--- ...
%     \item ... "--- ...
% \end{description}

% Раздел "Определения". Может отсутствовать в работе
% \definitions

% Раздел "Определения, обозначения и сокращения". Может отсутствовать в работе.
% Если присутствует, то заменяет собой разделы "Обозначения и сокращения" и
% "Определения"
% \defabbr

\section{Реализация}

\begin{Verbatim}[
  numbers=left,
  breaklines=true,
  breakanywhere=true
]
vector<int> buildBadCharTable(const string& pattern) {
    const int ALPHABET_SIZE = 256;
    vector<int> badCharTable(ALPHABET_SIZE, -1);

    for (int i = 0; i < pattern.size(); i++) {
        unsigned char currentChar = pattern[i];  //явно преобразуем в unsigned char
        int asciiCode = currentChar;              //получаем ascii код от 0 до 255
        badCharTable[asciiCode] = i;              //используем ascii код как индекс
    }

    return badCharTable;
}
\end{Verbatim}

\begin{Verbatim}[
  numbers=left,
  breaklines=true,
  breakanywhere=true
]
struct tree {
    int inf;
    tree* left;
    tree* right;
    int height;
};

tree* node(int x) {
    tree* n = new tree;
    n->inf = x;
    n->left = n->right = NULL;
    n->height = 1;
    return n;
}

int height(tree* t) {//считаем высоту
    if (t != nullptr) {
        return t->height;
    }
    else {
        return 0;
    }
}

int balance_factor(tree* t) {
    return height(t->right) - height(t->left);
}

void update_height(tree* t) {//обновление высоту
    int hl = height(t->left);
    int hr = height(t->right);

    if (hl > hr) {
        t->height = hl + 1;
    }
    else {
        t->height = hr + 1;
    }
}

tree* rotate_right(tree* t) {//правый поворот
    tree* q = t->left;
    t->left = q->right;
    q->right = t;
    update_height(t);
    update_height(q);
    return q;
}

tree* rotate_left(tree* t) {//левый поворот
    tree* p = t->right;
    t->right = p->left;
    p->left = t;
    update_height(t);
    update_height(p);
    return p;
}

tree* balance(tree* t) {//балансировка узла
    update_height(t);
    if (balance_factor(t) == 2) {
        if (balance_factor(t->right) < 0)
            t->right = rotate_right(t->right);
        return rotate_left(t);
    }
    if (balance_factor(t) == -2) {
        if (balance_factor(t->left) > 0)
            t->left = rotate_left(t->left);
        return rotate_right(t);
    }
    return t;
}

tree* insert(tree* t, int x) {//вставка нового элемента
    if (!t) return node(x);
    if (x < t->inf)
        t->left = insert(t->left, x);
    else
        t->right = insert(t->right, x);
    return balance(t);
}

tree* find_min(tree* t) {
    if (t->left != nullptr) {
        return find_min(t->left);
    }
    else {
        return t;
    }
}

tree* remove_min(tree* t) {//удаление минимального узла
    if (!t->left) return t->right;
    t->left = remove_min(t->left);
    return balance(t);
}

tree* remove(tree* t, int x) {//функция для удаления узла
    if (!t) return nullptr;

    if (x < t->inf) {
        t->left = remove(t->left, x);
    }
    else if (x > t->inf) {
        t->right = remove(t->right, x);
    }
    else {
        if (!t->left) {
            tree* temp = t->right;
            delete t;
            return temp;
        }
        else if (!t->right) {
            tree* temp = t->left;
            delete t;
            return temp;
        }
        else {
            tree* maxLeft = t->left;
            while (maxLeft->right) {
                maxLeft = maxLeft->right;
            }
            t->inf = maxLeft->inf;
            t->left = remove(t->left, maxLeft->inf);
        }
    }

    return balance(t);
}
\end{Verbatim}


\newpage

\section{Анализ}

АВЛ-дерево --- это самобалансирующееся двоичное дерево поиска, в котором для каждого узла выполняется свойство: абсолютная разница высот левого и правого поддеревьев (баланс-фактор) не превышает 1.
Пусть \(n\) --- количество узлов в дереве, \(h\) — высота дерева. Для АВЛ-дерева \(h = O(\log n)\) всегда.

\subsection*{Вставка элемента (insert)}

Функция \texttt{insert} реализована рекурсивно и выполняет вставку нового элемента с последующей балансировкой всех узлов на пути от листа к корню.

\textbf{Этап 1. Рекурсивный спуск к месту вставки:}
\begin{verbatim}
if (x < t->inf) t->left = insert(t->left, x);
else t->right = insert(t->right, x);
\end{verbatim}
Рекурсивный спуск от корня до листа требует \(h\) рекурсивных вызовов. Временная сложность спуска: \(O(h) = O(\log n)\).

\textbf{Этап 2. Создание нового узла:}
На дне рекурсии создаётся новый узел вызовом \texttt{node(x)} — \(O(1)\).

\textbf{Этап 3. Балансировка при возврате из рекурсии:}
После каждого рекурсивного вызова вызывается функция \texttt{balance(t)}:
\begin{itemize}
    \item Обновление высоты текущего узла через \texttt{update\_height} — \(O(1)\)
    \item Вычисление баланс-фактора — \(O(1)\)
    \item При необходимости выполнение поворотов (\texttt{rotate\_left}, \texttt{rotate\_right}) — каждый поворот выполняется за \(O(1)\)
\end{itemize}
Каждый вызов \texttt{balance} выполняется за \(O(1)\), и таких вызовов происходит не более \(h\) (на каждом уровне возврата из рекурсии). Однако важно отметить, что после выполнения поворота балансировка на вышележащих уровнях может не потребоваться, но в худшем случае балансировка выполняется на всех уровнях.

\textbf{Общая временная сложность вставки:}
\[
T_{\text{insert}}(n) = O(h) + O(1) + O(h) = O(\log n)
\]

\subsection*{Поиск элемента (find)}

Функция \texttt{find} выполняет стандартный поиск в двоичном дереве. На каждом шаге рекурсии происходит спуск на один уровень вниз. В худшем случае искомый элемент находится в листе или отсутствует — количество рекурсивных вызовов равно высоте дерева \(h\).

\[
T_{\text{find}}(n) = O(h) = O(\log n)
\]

\subsection*{Поиск минимума (find\_min)}

Функция \texttt{find\_min} рекурсивно спускается по левой ветви дерева до самого левого узла. Количество рекурсивных вызовов не превышает высоты дерева \(h\).

\[
T_{\text{find\_min}}(n) = O(h) = O(\log n)
\]

\subsection*{Удаление элемента (remove)}

Функция \texttt{remove} реализована рекурсивно и выполняет удаление элемента с последующей балансировкой.

\textbf{Этап 1. Рекурсивный спуск к удаляемому узлу:}
Аналогично поиску, требует \(O(h)\) рекурсивных вызовов: \(O(\log n)\).

\textbf{Этап 2. Обработка случаев при нахождении узла:}
\begin{itemize}
    \item \textbf{Нет левого потомка:} замена на правого потомка --- \(O(1)\)
    \item \textbf{Нет правого потомка:} замена на левого потомка --- \(O(1)\)
    \item \textbf{Два потомка:} поиск максимального элемента в левом поддереве требует спуска по правой ветви — \(O(h)\). Затем значение копируется, и рекурсивно удаляется найденный узел (имеющий не более одного потомка).
\end{itemize}

\textbf{Этап 3. Балансировка при возврате из рекурсии:}
После каждого рекурсивного вызова вызывается \texttt{balance(t)}. В отличие от вставки, при удалении балансировка может потребоваться на нескольких уровнях. В худшем случае балансировка выполняется на всех \(h\) уровнях возврата из рекурсии, каждый вызов \texttt{balance} — \(O(1)\).

\textbf{Общая временная сложность удаления:}
\[
T_{\text{remove}}(n) = O(h) + O(h) + O(h) = O(\log n)
\]

\subsection*{Удаление минимального элемента (remove\_min)}

Функция \texttt{remove\_min} рекурсивно спускается по левой ветви для удаления самого левого узла. Количество рекурсивных вызовов не превышает высоты дерева \(h\). В конце каждого вызова выполняется балансировка.

\[
T_{\text{remove\_min}}(n) = O(h) = O(\log n)
\]

\subsection*{Пространственная сложность (расход памяти)}

\textbf{Память для хранения дерева:}
Каждый узел структуры \texttt{tree} содержит:
\begin{itemize}
    \item \texttt{int inf} — 4 байта
    \item \texttt{tree* left} — указатель (4 или 8 байт)
    \item \texttt{tree* right} — указатель (4 или 8 байт)
    \item \texttt{int height} — 4 байта
\end{itemize}
Итого на один узел: \(8 + 2 \times \text{sizeof(pointer)}\) байт (на 64-битной системе — 24 байта). Для \(n\) узлов:
\[
M_{\text{storage}}(n) = O(n)
\]

\textbf{Дополнительная память на стеке вызовов:}
Рекурсивные функции (\texttt{insert}, \texttt{remove}, \texttt{find}, \texttt{find\_min}, \texttt{remove\_min}) используют стек вызовов глубиной до \(h = O(\log n)\).
\[
M_{\text{stack}}(n) = O(\log n)
\]

\textbf{Общая пространственная сложность:}
\[
M(n) = O(n) + O(\log n) = O(n)
\]

\end{document}
